%BeginMC
\question
Retinitis pigmentosus (RP) is a rare form of blindness, but has a high frequency on Tristen de Cu\~nha, a small island in the South Atlantic Ocean.  The island was colonized by 15 Britons during the 1800s.  One person carried one allele for RP.  Today, of the 240 people descended from the original 15, 4 are blind and 9 are carriers for the disease.  Which of the following best explains the high proportion of RP on the island?

\begin{choices}
\correctchoice founder effect and inbreeding (a form of non-random mating).%EndChoice
\choice founder effect and natural selection.%EndChoice
\choice genetic drift and mutation.%EndChoice
\choice bottleneck and inbreeding.%EndChoice
\choice inbreeding and mutation.%EndChoice
\end{choices}
%EndMC

%BeginMC
\question
Prairie chickens are endangered grassland birds. In Illinois, 
grasslands used to cover most of the state but farming and growth of 
cities changed the habitat over a fairly short period of time. Now, 
less than 1\% of the grasslands remain. A study showed that most genes 
had far fewer alleles than historic populations. The loss of genetic 
variation is most likely due to

\begin{choices}
\choice the small population size causing random fluctuation of allele frequencies. %EndChoice
\correctchoice a genetic bottleneck caused by the loss of habitat. %EndChoice
\choice a founder event caused by individuals moving to the remaining grasslands. %EndChoice
\choice directional selection that favored survivorship in small habitats. %EndChoice
\choice balancing selection, which maintains genetic variation in small populations. %EndChoice
\end{choices}
%EndMC

%BeginMC
\question
In science, a hypothesis is only useful if it

\begin{choices}
	\correctchoice can be tested. %EndChoice
	\choice can be proven correct. %EndChoice
	\choice is proven correct. %EndChoice
	\choice explains the natural world. %EndChoice
	\choice agrees with the conclusion. %EndChoice
\end{choices}
%EndMC

%BeginMC
\question
The males of some fishes and birds have bright colors that males use during territorial displays. The existence of these bright colors in unrelated organisms is an example of 

\begin{choices}
\choice homology. %EndChoice
\choice common ancestry. %EndChoice
\correctchoice convergent evolution. %EndChoice
\choice descent with modification. %EndChoice
\choice intersexual selection. %EndChoice
\end{choices}
%EndMC


%BeginMC
\question
\textit{Anopheles} mosquitoes, which carry the malaria parasite, cannot live above elevations of 5,900 feet. In addition, oxygen availability decreases with higher altitude. Consider a hypothetical human population that is adapted to life on the slopes of Mt. Kilimanjaro in Tanzania, a country in equatorial Africa. Mt. Kilimanjaro's base is about 2,600 feet above sea level and its peak is 19,341 feet above sea level. If the frequency of the sickle-cell allele in the population is plotted against altitude (feet above sea level), which of the following distributions is most likely, assuming little migration of people up or down the mountain?

\begin{multicols}{2}
\begin{choices}
\choice \includegraphics[width=0.35\textwidth]{anopA} %EndChoice
\correctchoice \includegraphics[width=0.35\textwidth]{anopB} %EndChoice
\choice \includegraphics[width=0.35\textwidth]{anopC} %EndChoice
\choice \includegraphics[width=0.35\textwidth]{anopD} %EndChoice
\end{choices}
\end{multicols}
%EndMC

%BeginMC
\question
Over time, the movement of people among all countries on Earth has steadily increased. This has altered the course of human evolution by increasing

\begin{choices}
\choice nonrandom mating. %EndChoice
\choice geographic isolation. %EndChoice
\choice genetic drift. %EndChoice
\correctchoice gene flow. %EndChoice
\choice natural selection. %EndChoice
\end{choices}
%EndMC


%BeginMC
\question
Which two violations of the Hardy-Weinberg assumptions \emph{always} causes homozygosity to increase in a population?

\begin{choices}
\correctchoice Genetic drift and non-random mating.%EndChoice
\choice Mutation and natural selection.%EndChoice
\choice Natural selection and gene flow.%EndChoice
\choice Genetic drift and natural selection.%EndChoice
\choice Non-random mating and gene flow.%EndChoice
\choice Gene flow and mutation.%EndChoice
\end{choices}
%EndMC



%BeginMC
\question
A population in Hardy-Weinberg equilibrium has alleles \textit{$A_1$} and \textit{$A_2$}. The frequency of allele \textit{$A_2$} is 0.2. What is the frequency of heterozygotes? 

\begin{choices}
	\choice 0.4 %dontshuffle%EndChoice
	\choice 0.02 %dontshuffle%%EndChoice
	\choice 0.04 %dontshuffle%%EndChoice
	\choice 0.16 %dontshuffle%%EndChoice
	\correctchoice 0.32 %dontshuffle%%EndChoice
	\choice Not enough information to determine. %dontshuffle%%EndChoice
\end{choices}
%EndMC

%BeginMC
\question
In a population with two alleles, $A_1$ and $A_2$, that is in Hardy-Weinberg equilibrium, the frequency of allele $A_1$ is 0.2. What is the frequency of $A_2A_2$ individuals? 

\begin{choices}
	\choice 0.4 %dontshuffle%EndChoice
	\choice 0.02 %dontshuffle%EndChoice
	\choice 0.04 %dontshuffle%EndChoice
	\choice 0.16 %dontshuffle%EndChoice
	\choice 0.32 %dontshuffle%EndChoice
	\correctchoice 0.64 %dontshuffle%EndChoice
	\choice Not enough information to determine. %dontshuffle%%EndChoice
\end{choices}
%EndMC

%BeginMC
\question
A species of sport fish reaches an average adult length of 100 cm. Fishing regulations require that only individuals longer than 75 cm may be kept. Shorter fish must be released because they have greater reproductive output relative to large individuals. Based on your knowledge of natural selection, you should predict that over hundreds of generations the average length of the adult fish population will

\begin{choices}
\correctchoice remain unchanged. %EndChoice
\choice gradually decline. %EndChoice
\choice rapidly decline. %EndChoice
\choice gradually increase. %EndChoice
\choice rapidly increase. %EndChoice
\end{choices}
%EndMC


%BeginMC
\question
The rabbitsfoot mussel was once widespread throughout rivers of eastern North America but now remains in only a few very small populations. The decrease occurred rapidly due to habitat loss. Genetics studies show that the remaining populations have little or no genetic variation. The best explanation for this result is:

\begin{choices}
\choice Stabilizing selection favored only one or two genes. %EndChoice
\choice Genetic drift caused the loss of homozygosity. %EndChoice
\choice Inbreeding caused a decrease in homozygosity. %EndChoice
\correctchoice A bottleneck caused the increase of homozygosity. %EndChoice
\choice A founder event caused the loss of heterozygosity. %EndChoice
\end{choices}
%EndMC

%BeginMC
\question
The rabbitsfoot mussel was once widespread throughout rivers of eastern North America but now remains in only a few very small populations. The decrease occurred rapidly due to habitat loss. Genetics studies show that the remaining populations have little or no genetic variation. The best explanation for this result is:

\begin{choices}
\choice Balancing selection favored only one or two alleles. %EndChoice
\correctchoice Genetic drift caused the loss of heterozygosity. %EndChoice
\choice Inbreeding caused an increase in heterozygosity. %EndChoice
\choice Directional selection favored only one or two alleles. %EndChoice
\choice A bottleneck caused the loss of homozygosity. %EndChoice
\end{choices}
%EndMC

%BeginMC
\question
The island of Krakatou is a volcanic island in Indonesia.  An eruption in 1883 killed all living organisms on the island.  Many of organismal populations that are present now tend to show low genetic variation.  This is likely due to

\begin{choices}
\choice directional selection that favored traits that allow organisms to withstand volcanic activity. %EndChoice
\choice balancing selection that favored traits that allow organisms to withstand volcanic activity. %EndChoice
\choice a genetic bottleneck caused by the eruption, followed by regrowth of the population regrew. %EndChoice
\correctchoice a series of founder effects with different species recolonizing from surrounding islands. %EndChoice
\choice The small population size causing random fluctuation of allele frequencies. %EndChoice
\end{choices}
%EndMC


%BeginMC
\question
Retinitis pigmentosus (RP) is a rare form of blindness, but has a high frequency on Tristen de Cu\~nha, a small island in the South Atlantic Ocean.  The island was colonized by 15 Britons during the 1800s.  One person carried one allele for RP.  Today, of the 240 people descended from the original 15, 4 are blind and 9 are carriers for the disease.  Which of the following best explains the high proportion of RP on the island?

\begin{choices}
\correctchoice founder effect and inbreeding (a form of non-random mating).%EndChoice
\choice founder effect and natural selection.%EndChoice
\choice genetic drift and mutation.%EndChoice
\choice bottleneck and inbreeding.%EndChoice
\choice inbreeding and mutation.%EndChoice
\end{choices}
%EndMC


%BeginMC
\question
Some beetles and flies have antler-like structures on their heads, much like male deer do. The existence of antlers in beetle, fly, and deer species with strong male-male competition is an example of 

\begin{choices}
\correctchoice convergent evolution. %EndChoice
\choice intersexual selection. %EndChoice
\choice homology. %EndChoice
\choice common ancestry. %EndChoice
\end{choices}
%EndMC


%BeginMC
\question
Which one of the following most accurately describes the process by which scientific knowledge advances? 

\begin{choices}
\choice The use of scientific analysis to explain the natural world by theory and experiments.%EndChoice
\correctchoice The objective analysis of the natural world by hypothesis testing and experiments.%EndChoice
\choice The analysis of the natural world through theoretical reasoning and experiments.%EndChoice
\choice The use of deductive reasoning to develop broad scientific principles that explain the natural world.%EndChoice
\choice The use of inductive reasoning to explain specific experimental results about the natural world.%EndChoice
\end{choices}
%EndMC


%BeginMC
\question
In science, a hypothesis is only useful if it

\begin{choices}
\correctchoice can be tested.%EndChoice
\choice can be proven correct.%EndChoice
\choice is proven correct.%EndChoice
\choice explains the natural world.%EndChoice
\choice agrees with the conclusion.%EndChoice
\end{choices}
%EndMC

%BeginMC
\question
After invention of the microscope, scientists discovered that all living organisms were composed of one or more cells, that the cell was the basic unit of function in all organisms, and that only existing cells gave rise to new cells. These specific observations were eventually organized into a general theory of biology called the Cell Theory. This is an example of 

\begin{choices}
\choice deductive reasoning.%EndChoice
\choice hypothetical reasoning.%EndChoice
\correctchoice inductive reasoning.%EndChoice
\choice objective reasoning.%EndChoice
\choice scientific reasoning.%EndChoice
\choice theoretical reasoning.%EndChoice
\end{choices}
%EndMC

%BeginMC
\question
Which mouse has the greatest relative fitness? (A clutch is the offspring from one breeding event.)

\begin{choices}
\choice One that lives for 4 years, has 3 clutches per year, and 4 offspring per clutch.%EndChoice
\choice One that lives for 3 years, has 3 clutches per year, and 5 offspring per clutch.%EndChoice 
\choice One that lives for 3 years, has 2 clutches per year, and 6 offspring per clutch. %EndChoice
\choice One that lives for 2 years, has 3 clutches per year, and 8 offspring per clutch. %EndChoice
\correctchoice One that lives for 3 years, has 3 clutches per year, and 7 offspring per clutch. %EndChoice
\end{choices}
%EndMC

%BeginMC
\question
Which of these conditions should completely prevent the occurrence of natural selection in a population over time?

\begin{choices}
\correctchoice All variation among individuals is due only to environmental factors. %EndChoice
\choice The environment is changing at a relatively slow rate. %EndChoice
\choice The population size is large. %EndChoice
\choice The population lives in a habitat where there are no competing species present. %EndChoice
\end{choices}
%EndMC

%BeginMC
\question
According to Darwin's concept of descent with modification:

\begin{choices}
\choice Populations evolve by natural selection acting on individuals. %EndChoice
\correctchoice Species evolve from ancestral forms through the accumulation of different adaptations. %EndChoice
\choice Individuals with greater fitness will evolve more quickly than individuals with lower fitness. %EndChoice
\choice All organisms overproduce offspring, leading to differences in survivorship. %EndChoice
\choice All organisms overproduce offspring, leading to different adaptations. %EndChoice
\end{choices}
%EndMC

% This choice used to be part of the answer above. Removed to reduce some confusion.
\choice All species have evolved from a single common ancestor. %EndChoice


%BeginMC
\question
The overuse of drugs such as penicillin has caused an increase in the number of bacterial species that are resistant to the drugs. The most likely explanation for this result is:

\begin{choices}
\choice Most bacteria can resist drugs. %EndChoice
\choice Penicillin was not an effective drug. %EndChoice
\choice Bacteria began making proteins resistant to the drugs. %EndChoice
\choice The drugs caused mutations in the bacteria that gave them resistance. %EndChoice
\correctchoice A few bacteria were naturally resistant and natural selection favored their increase. %EndChoice
\end{choices}
%EndMC

%BeginMC
\question
Organisms living in the northern rocky intertidal face competing selective forces. Larger individuals benefit by having a larger internal volume, which helps to conserve body temperature. However, smaller individuals are less likely to be dislodged by large waves from storms. Thus, selection may favor intermediate body size for any particular species. This is an example of

\begin{choices}
\choice directional selection. %EndChoice
\correctchoice stabilizing selection. %EndChoice
\choice disruptive selection. %EndChoice
\choice balancing selection. %EndChoice
\choice natural selection. %EndChoice
\end{choices}
%EndMC

%BeginMC
\question
Snails living in the southern rocky intertidal have adapted to warm temperatures. Compared to northern snails, southern snails have ridges to get rid of excess heat and lighter colors to reflect heat. If the snail phenotype ranges from no ridges and dark color to ridges and light color, then the mode of selection that fits the southern snails is most likely

\begin{choices}
	\correctchoice directional selection. %EndChoice
	\choice stabilizing selection. %EndChoice
	\choice disruptive selection. %EndChoice
	\choice sexual selection. %EndChoice
	\choice balancing selection. %EndChoice
\end{choices}
%EndMC


%BeginMC
\question
Which mode of selection best applies to this example: female widowbirds choose males with the longest tails for mating. However, males with very long tails are caught more often by predators because they are slower to take flight. A long-term study showed that male tail length is not getting longer or shorter. 

\begin{choices}
	\choice directional selection. %EndChoice
	\correctchoice stabilizing selection. %EndChoice
	\choice disruptive selection. %EndChoice
	\choice sexual selection. %EndChoice
	\choice balancing selection. %EndChoice
\end{choices}
%EndMC



%BeginMC
\question
Which of the following is (are) true of natural selection (choose only one)?

\begin{choices}
\choice It requires genetic variation. %EndChoice
\choice It results in descent with modification. %EndChoice
\choice It involves differences in reproductive success. %EndChoice
\choice It results in descent with modification and involves differences in reproductive success. %EndChoice
\correctchoice It requires genetic variation, results in descent with modification, and involves differences in reproductive success. %EndChoice
\end{choices}
%EndMC

%BeginMC
\question
Given a population that contains genetic variation, what is the correct sequence of the following events, under the influence of natural selection?
\begin{enumerate}
	\item Well-adapted individuals leave more offspring than do poorly adapted individuals.
	\item A change occurs in the environment.
	\item Genetic frequencies within the population change.
	\item Poorly adapted individuals have decreased survivorship.
\end{enumerate}

\begin{choices}
\correctchoice 2 $\rightarrow$ 4 $\rightarrow$ 1 $\rightarrow$ 3 %EndChoice
\choice 4 $\rightarrow$ 2 $\rightarrow$ 1 $\rightarrow$ 3 %EndChoice
\choice 4 $\rightarrow$ 1 $\rightarrow$ 2 $\rightarrow$ 3 %EndChoice
\choice 4 $\rightarrow$ 2 $\rightarrow$ 3 $\rightarrow$ 1 %EndChoice
\choice 2 $\rightarrow$ 4 $\rightarrow$ 3 $\rightarrow$ 1 %EndChoice
\end{choices}
%EndMC


%%%=====NO LONGER USED.


%BeginMC
\question
According to your readings, one of the common themes of life is that life requires the transfer and transformation of energy and matter.  In living organisms, this occurs through

\begin{choices}
\choice homeostasis. %EndChoice
\correctchoice metabolism. %EndChoice
\choice growth and development. %EndChoice
\choice reproduction and heredity. %EndChoice
\choice complexity and organization. %EndChoice
\choice response to the environment. %EndChoice
\end{choices}
%EndMC

%BeginMC
\question
(Think carefully about this.) Organisms living in the northern rocky intertidal face competing selective forces. Larger individuals benefit by having a larger internal volume, which helps to conserve body temperature. However, smaller individuals are less likely to be dislodged by large waves from storms. Thus, selection may favor intermediate body size for any particular species. This is an example of

\begin{choices}
\choice directional selection. %EndChoice
\choice stabilizing selection. %EndChoice
\choice disruptive selection. %EndChoice
\correctchoice balancing selection. %EndChoice
\choice natural selection. %EndChoice
\end{choices}
%EndMC

%BeginMC
\question
According to your text, which of the following represents the biological hierarchy from organism to biosphere:

\begin{choices}
\correctchoice Organism: Population: Community: Ecosystem: Biosphere %EndChoice
\choice Organism: Community: Population: Ecosystem: Biosphere %EndChoice
\choice Organism: Ecosystem: Community: Population: Biosphere %EndChoice
\choice Organism: Population: Ecosystem: Community: Biosphere %EndChoice
\choice Organism: Community: Ecosystem: Population: Biosphere %EndChoice
\end{choices}
%EndMC


%BeginMC
\question
Which one of the following is the best definition of a hypothesis?

\begin{choices}
\choice An alternative explanation for a null hypothesis.%EndChoice
\choice A tentative explanation about a supernatural phenomenon.%EndChoice
\choice A brief summary of a broad scientific theory.%EndChoice
\correctchoice A tentative explanation for a natural phenomenon.%EndChoice
\choice A prediction about the expected results from a specific experiment.%EndChoice
\end{choices}
%EndMC

%BeginMC
\question
Which one of the following statements \emph{cannot} be tested by the scientific method?

\begin{choices}
\choice The order in which organisms, including humans, first appeared on earth. For example, man, then plants, then land animals and birds, then woman.%EndChoice
\choice The Grand Canyon was formed by a global flood about 4000 years ago.%EndChoice
\choice The Earth is between 6000 and 10,000 years old.%EndChoice
\correctchoice God created light and dark, then stars.%EndChoice
\choice The entire universe, the earth, and all living things were created in six days.%EndChoice
\choice All of these statements are religious claims so none can be tested by the scientific method.%dontshuffle%EndChoice
\end{choices}
%EndMC
